\documentclass[11pt, a4paper]{article}

% --- Packages ---
\usepackage[a4paper, top=2.5cm, bottom=2.5cm, left=2.5cm, right=2.5cm]{geometry}
\usepackage{amsmath, amssymb}
\usepackage{graphicx}
\usepackage{booktabs}
\usepackage{parskip}
\usepackage{microtype}
\usepackage{fancyhdr}
\usepackage{titlesec}
\usepackage{xcolor}
\usepackage{mdframed}
\usepackage{enumitem}
\usepackage{hyperref}
\usepackage{caption}
\usepackage{float}
\usepackage{array}
\usepackage{multirow}
\usepackage{colortbl}

\hypersetup{colorlinks=true, linkcolor=black, urlcolor=blue, citecolor=black}

% --- Page style ---
\pagestyle{fancy}
\fancyhf{}
\renewcommand{\headrulewidth}{0.4pt}
\renewcommand{\footrulewidth}{0.4pt}
\fancyhead[L]{\small\textit{Study Notes: C-MAPSS Dataset — Deep Reference}}
\fancyhead[R]{\small\thepage}
\fancyfoot[C]{\small\textit{``I'm not stupid, I'm just slow''}}

% --- Section style ---
\titleformat{\section}{\large\bfseries}{}{0em}{}[\titlerule]
\titleformat{\subsection}{\normalsize\bfseries}{}{0em}{}
\titleformat{\subsubsection}{\normalsize\itshape}{}{0em}{}

% --- Highlight boxes ---
\newmdenv[
  backgroundcolor=gray!8,
  linecolor=gray!40,
  linewidth=0.8pt,
  innerleftmargin=10pt,
  innerrightmargin=10pt,
  innertopmargin=8pt,
  innerbottommargin=8pt,
  skipabove=8pt,
  skipbelow=8pt
]{keybox}

\newmdenv[
  backgroundcolor=orange!6,
  linecolor=orange!40,
  linewidth=0.8pt,
  innerleftmargin=10pt,
  innerrightmargin=10pt,
  innertopmargin=8pt,
  innerbottommargin=8pt,
  skipabove=8pt,
  skipbelow=8pt
]{insightbox}

% --- Table row colors ---
\definecolor{grpA}{RGB}{225, 245, 238}
\definecolor{grpB}{RGB}{238, 237, 254}
\definecolor{grpC}{RGB}{250, 238, 218}
\definecolor{grpD}{RGB}{250, 236, 231}
\definecolor{grpE}{RGB}{230, 241, 251}
\definecolor{grpF}{RGB}{234, 243, 222}
\definecolor{grpG}{RGB}{251, 234, 240}
\definecolor{grpH}{RGB}{241, 239, 232}

\captionsetup{font=small, labelfont=bf, skip=4pt}

% -------------------------------------------------------------------
\begin{document}
% -------------------------------------------------------------------

{\Large\bfseries Dataset Notes: C-MAPSS --- Deep Reference}

\smallskip
{\small Saxena, Goebel, Simon, Eklund \quad --- \quad PHM'08 dataset \quad --- \quad
\textit{Column-by-column edition}}

\medskip\hrule

\tableofcontents

\bigskip\hrule\bigskip

% -------------------------------------------------------------------
\section{Overview: Four Sub-Datasets}
% -------------------------------------------------------------------

Figure~\ref{fig:matrix} shows the four C-MAPSS sub-datasets and their
key properties. All share the same 26-column file format.

\begin{figure}[H]
  \centering
  \includegraphics[width=\textwidth]{09_dataset_matrix.png}
  \caption{The four C-MAPSS sub-datasets. FD001 is simplest
  (1 condition, 1 fault mode). FD004 is hardest
  (6 conditions, 2 fault modes).}
  \label{fig:matrix}
\end{figure}

% -------------------------------------------------------------------
\section{The Three Files and Their Different Lengths}
% -------------------------------------------------------------------

\subsection{Why the lengths differ}

Each sub-dataset produces three files. They differ in length because
they represent the same engines at different levels of granularity.

\begin{figure}[H]
  \centering
  \includegraphics[width=\textwidth]{18_sheet_length_relationship.png}
  \caption{The three sheets for FD001. Train and test are
  cycle-level (many rows per engine). RUL is engine-level
  (one row per engine). The bottom table summarises the
  three granularities.}
  \label{fig:lengths}
\end{figure}

\begin{center}
\small
\begin{tabular}{@{}llrp{6cm}@{}}
  \toprule
  File & Granularity & FD001 rows & Why \\
  \midrule
  \texttt{train\_FD001.txt} & one row per cycle & 20,631
    & 100 engines $\times$ avg 206 cycles to failure \\
  \texttt{test\_FD001.txt}  & one row per cycle & 13,096
    & 100 engines $\times$ avg 131 cycles (truncated) \\
  \texttt{RUL\_FD001.txt}   & one row per engine & 100
    & one RUL value per test engine \\
  \bottomrule
\end{tabular}
\end{center}

\subsection{How the three files connect}

The link between all three files is \texttt{unit\_number}.
It is the key for any join operation.

\begin{figure}[H]
  \centering
  \includegraphics[width=0.92\textwidth]{10_train_test_rul_relationship.png}
  \caption{The relationship between the three files for one engine.
  The train file runs to failure. The test file stops at $t^*$.
  The RUL file gives one number per engine: the true cycles
  remaining at $t^*$.}
  \label{fig:filerel}
\end{figure}

\begin{keybox}
\textbf{The three files are three zoom levels of the same engines:}

\medskip
\begin{tabular}{@{}lll@{}}
  \texttt{train} & cycle-level & full life, all sensors, all cycles \\
  \texttt{test}  & cycle-level & partial life, sensors only \\
  \texttt{RUL}   & engine-level & one number: true RUL at cut point \\
\end{tabular}

\medskip
They are connected by \texttt{unit\_number}. Analysis that combines
files must always join on this column.
\end{keybox}

\subsection{What this means for analysis}

Different questions require different combinations:

\begin{center}
\small
\begin{tabular}{@{}lp{4.5cm}p{5cm}@{}}
  \toprule
  Question & Files used & How \\
  \midrule
  What does degradation look like? &
    train only &
    Plot sensor vs cycle for each engine \\
  How well did I predict RUL? &
    test + RUL, join on unit\_number &
    Compare predicted RUL to true RUL \\
  What is the engine lifespan distribution? &
    RUL + last\_cycle (v3) &
    total\_life\_est = last\_cycle + RUL \\
  Do long-lived engines degrade differently? &
    train + RUL, join on unit\_number &
    Group train rows by engine life length \\
  \bottomrule
\end{tabular}
\end{center}

% -------------------------------------------------------------------
\section{Input Columns: All 26 Raw Columns}
% -------------------------------------------------------------------

Every \texttt{train} and \texttt{test} row has 26 columns.
They are grouped below by physical meaning.
Figure~\ref{fig:engine} shows where each group lives on the engine.

\begin{figure}[H]
  \centering
  \includegraphics[width=\textwidth]{03_cmapss_engine_schematic.png}
  \caption{C-MAPSS turbofan engine. Sensors measure temperatures
  and pressures at inlet, inter-stage, and outlet locations.
  The amber HPC module is the one degraded in FD001/FD002.}
  \label{fig:engine}
\end{figure}

\begin{table}[H]
\centering
\caption{All 26 raw input columns grouped by physical meaning.}
\label{tab:allcols}
\small
\begin{tabular}{@{}r >{\ttfamily}l l p{6.8cm}@{}}
\toprule
Col & \normalfont Name & Units & Physical meaning \\
\midrule

\rowcolor{grpA}
\multicolumn{4}{@{}l}{\textbf{Group A — Identifiers}} \\
\rowcolor{grpA}
1  & unit\_number  & ---    &
  Engine ID. Integer. Resets per file.
  Every analysis groups or filters on this. \\
\rowcolor{grpA}
2  & time\_cycles  & cycles &
  Cycle counter per engine. Starts at 1.
  One cycle $\approx$ one flight. \\

\rowcolor{grpB}
\multicolumn{4}{@{}l}{\textbf{Group B — Operating conditions}} \\
\rowcolor{grpB}
3  & op\_setting\_1 & ft    &
  Altitude. Controls ambient pressure and temperature.
  Ranges 0--42,000\,ft across 6 conditions. \\
\rowcolor{grpB}
4  & op\_setting\_2 & ---   &
  Mach number. Flight speed as fraction of speed of sound.
  Ranges 0--0.84. \\
\rowcolor{grpB}
5  & op\_setting\_3 & ---   &
  Throttle resolver angle (TRA). Commanded thrust level.
  Ranges 20--100. Constant within FD001/FD003. \\

\rowcolor{grpC}
\multicolumn{4}{@{}l}{\textbf{Group C — Temperatures} (total = stagnation, measures enthalpy)} \\
\rowcolor{grpC}
6  & T2     & \textdegree R &
  Total temperature at fan inlet. Reflects ambient conditions.
  Jumps dramatically across flight conditions. \\
\rowcolor{grpC}
7  & T24    & \textdegree R &
  Total temperature at LPC outlet / HPC inlet.
  Rises as air is compressed in LPC. \\
\rowcolor{grpC}
8  & T30    & \textdegree R &
  Total temperature at HPC outlet / combustor inlet.
  Key health indicator: rises when HPC degrades. \\
\rowcolor{grpC}
9  & T50    & \textdegree R &
  Total temperature at LPT outlet.
  Reflects turbine work extraction and combustor output. \\

\rowcolor{grpD}
\multicolumn{4}{@{}l}{\textbf{Group D — Pressures}} \\
\rowcolor{grpD}
10 & P2     & psia &
  Total pressure at fan inlet. Ambient total pressure.
  Drops at altitude. \\
\rowcolor{grpD}
11 & P15    & psia &
  Total pressure in bypass duct (between fan and LPC).
  Reflects fan performance. \\
\rowcolor{grpD}
12 & P30    & psia &
  Total pressure at HPC outlet. Main compression output.
  Falls as HPC degrades. \\
\rowcolor{grpD}
16 & Ps30   & psia &
  Static pressure at HPC outlet (no velocity component).
  Related to P30 but measured differently. \\

\rowcolor{grpE}
\multicolumn{4}{@{}l}{\textbf{Group E — Speeds}} \\
\rowcolor{grpE}
13 & Nf     & rpm &
  Physical fan speed. Actual rotational speed of the fan shaft.
  Controlled by fan-speed controller. \\
\rowcolor{grpE}
14 & Nc     & rpm &
  Physical core speed. Actual rotational speed of the
  HPC/HPT/combustor shaft. \\
\rowcolor{grpE}
18 & NRf    & rpm &
  Corrected fan speed. Adjusted for inlet conditions:
  $NR_f = N_f / \sqrt{T_2/T_{\text{ref}}}$.
  Comparable across conditions. \\
\rowcolor{grpE}
19 & NRc    & rpm &
  Corrected core speed. Same correction applied to $N_c$:
  $NR_c = N_c / \sqrt{T_2/T_{\text{ref}}}$.
  Comparable across conditions. \\

\rowcolor{grpF}
\multicolumn{4}{@{}l}{\textbf{Group F — Flow ratios and combustion}} \\
\rowcolor{grpF}
15 & epr    & --- &
  Engine pressure ratio $P_{50}/P_2$. Overall pressure
  rise through the engine. Thrust indicator. \\
\rowcolor{grpF}
17 & phi    & pps/psi &
  Fuel flow / $P_{s30}$ ratio. Normalised fuel flow.
  Rises when HPC pressure drops due to degradation. \\
\rowcolor{grpF}
20 & BPR    & --- &
  Bypass ratio. Mass flow through bypass duct divided by
  core mass flow. Reflects fan/core balance. \\
\rowcolor{grpF}
21 & farB   & --- &
  Burner fuel-air ratio. Mass of fuel per mass of air
  in the combustor. Combustion efficiency indicator. \\
\rowcolor{grpF}
22 & htBleed & --- &
  Bleed enthalpy. Energy extracted from the compressor
  for aircraft systems (cabin air, de-icing). \\

\rowcolor{grpG}
\multicolumn{4}{@{}l}{\textbf{Group G — Controller demands}} \\
\rowcolor{grpG}
23 & Nf\_dmd      & rpm &
  Demanded fan speed. The speed the controller is
  targeting. Compare to $N_f$ to see control error. \\
\rowcolor{grpG}
24 & PCNfR\_dmd   & rpm &
  Demanded corrected fan speed.
  The corrected-speed target from the controller. \\

\rowcolor{grpH}
\multicolumn{4}{@{}l}{\textbf{Group H — Coolant bleeds}} \\
\rowcolor{grpH}
25 & W31    & lbm/s &
  HPT coolant bleed. Air bled from HPC to cool the
  high-pressure turbine blades. \\
\rowcolor{grpH}
26 & W32    & lbm/s &
  LPT coolant bleed. Air bled from HPC to cool the
  low-pressure turbine blades. \\
\bottomrule
\end{tabular}
\end{table}

\begin{insightbox}
\textbf{Why corrected speeds (NRf, NRc) matter for FD002/FD004.}
Physical speeds $N_f$ and $N_c$ depend on inlet temperature.
At high altitude (cold air), the same thrust requires different RPM
than at sea level. Corrected speeds remove this dependency, making
the values comparable across flight conditions. For single-condition
sub-datasets (FD001/FD003) this distinction is less important.
\end{insightbox}

\begin{keybox}
\textbf{Which columns carry degradation signal?}
HPC degradation (FD001/FD002) primarily shows up in:
T30 $\uparrow$ (outlet temperature rises),
P30 $\downarrow$ (outlet pressure falls),
phi $\uparrow$ (more fuel needed to maintain thrust),
Ps30 $\downarrow$.
The other sensors change too, but these four are the most direct.
\end{keybox}

% -------------------------------------------------------------------
\section{Additional Input Columns}
% -------------------------------------------------------------------

These three columns are not in the raw files. They are derived
from the 26 raw columns and appended in the v3 spreadsheet.
They appear on both \texttt{train} and \texttt{test} sheets.

\subsection{condition\_id}

For FD002 and FD004, raw sensor values are not comparable across
cycles because the engine operates at different flight conditions.
Figure~\ref{fig:condproblem} shows this problem clearly.

\begin{figure}[H]
  \centering
  \includegraphics[width=\textwidth]{12_multi_condition_problem.png}
  \caption{Top: raw $T_2$ values jumping between six bands as the
  engine switches flight conditions --- not degradation signal.
  Bottom: after condition-wise normalisation, the health signal
  becomes visible.}
  \label{fig:condproblem}
\end{figure}

\texttt{condition\_id} assigns an integer 1--6 to each row
identifying which of the six flight conditions that cycle belongs to.
It is computed by KMeans clustering on the three operating setting
columns, as shown in Figure~\ref{fig:normpipeline}.

\begin{figure}[H]
  \centering
  \includegraphics[width=0.82\textwidth]{13_normalization_pipeline.png}
  \caption{How \texttt{condition\_id} is produced.
  KMeans$(k=6)$ is fitted on the train set and applied to both
  train and test. For FD001/FD003, all rows receive
  \texttt{condition\_id} = 1.}
  \label{fig:normpipeline}
\end{figure}

The KMeans model minimises within-cluster sum of squared distances:

\begin{equation}
  J = \sum_{i=1}^{n} \left\| x_i - \mu_{c(i)} \right\|^2
  \label{eq:kmeans_obj}
\end{equation}

where $x_i = (\texttt{op\_setting\_1}_i,\, \texttt{op\_setting\_2}_i,\,
\texttt{op\_setting\_3}_i)$,
$\mu_k$ is the centroid of cluster $k$, and
$c(i) = \arg\min_k \|x_i - \mu_k\|^2$ is the assigned cluster.

\begin{keybox}
\textbf{Important:} \texttt{condition\_id} labels are not comparable
across sub-datasets. Condition 2 in FD002 and condition 2 in FD004
refer to different physical flight points because KMeans assigns
labels in arbitrary order. Always compare conditions by their
\texttt{op\_setting} values, not by label number.
\end{keybox}

\subsection{dT\_compressor}

Temperature rise across the HPC:

\begin{equation}
  \texttt{dT\_compressor} = T_{30} - T_{24}
  \label{eq:dtcomp}
\end{equation}

where $T_{30}$ is the total temperature at the HPC outlet (col 8)
and $T_{24}$ is the total temperature at the LPC outlet / HPC inlet
(col 7).

\textbf{Physical meaning:} healthy compression raises temperature
by a predictable amount. As the HPC degrades, compression efficiency
drops. The controller compensates by demanding more fuel, which raises
$T_{30}$ further. So \texttt{dT\_compressor} tends to
\emph{increase} as HPC health deteriorates.

\textbf{Why add this column?} The raw $T_{30}$ value is dominated
by the operating condition (altitude, throttle). The difference
$T_{30} - T_{24}$ removes the common-mode variation and isolates
the HPC contribution.

\subsection{dT\_turbine}

Temperature drop across the turbine section:

\begin{equation}
  \texttt{dT\_turbine} = T_{50} - T_{30}
  \label{eq:dtturb}
\end{equation}

where $T_{50}$ is the total temperature at the LPT outlet (col 9)
and $T_{30}$ is the HPC outlet / combustor inlet (col 8).

\textbf{Physical meaning:} the turbines extract work from the hot
combustion gases, dropping their temperature. A healthy turbine
produces a large drop. As turbine efficiency changes, this delta
shifts. Note: in FD001/FD002 the turbines are not the degraded
component, so this column is less discriminative there but remains
useful as a health indicator for FD003/FD004 where Fan degradation
can affect the thermodynamic balance.

% -------------------------------------------------------------------
\section{Output Columns (Train Sheet Only)}
% -------------------------------------------------------------------

These three columns are the \emph{targets} for machine learning.
They appear only on the \texttt{train} sheet because the test RUL
is the unknown quantity to be predicted.

Figure~\ref{fig:derivation} shows the full derivation map for all
additional columns.

\begin{figure}[H]
  \centering
  \includegraphics[width=\textwidth]{19_column_derivation.png}
  \caption{Derivation map for all additional columns.
  Source columns (gray, top row) feed into three groups
  of derived columns via different computation paths.}
  \label{fig:derivation}
\end{figure}

\subsection{RUL\_raw}

Remaining useful life in cycles, counting down to zero at failure:

\begin{equation}
  \texttt{RUL\_raw}(i) = t_{\max}(u_i) - t_i
  \label{eq:rul_raw}
\end{equation}

where:
\begin{itemize}[noitemsep]
  \item $t_i$ is \texttt{time\_cycles} for row $i$
  \item $u_i$ is \texttt{unit\_number} for row $i$
  \item $t_{\max}(u_i) = \max\{t_j : u_j = u_i\}$ is the last
        cycle observed for that engine (the failure cycle in the
        training set)
\end{itemize}

\texttt{RUL\_raw} = 0 at the last row of each training engine
and counts upward toward earlier cycles.

\textbf{Spreadsheet formula (LO Calc):}
\begin{verbatim}
{=MAX(IF($A:$A=$A2,$B:$B))-B2}
Ctrl+Shift+Enter (array formula)
\end{verbatim}

\subsection{RUL\_capped}

A piece-wise linear target that clamps long-range RUL values:

\begin{equation}
  \texttt{RUL\_capped}(i) = \min\!\bigl(\texttt{RUL\_raw}(i),\; 125\bigr)
  \label{eq:rul_capped}
\end{equation}

Figure~\ref{fig:rulshape} shows the difference between the two shapes.

\begin{figure}[H]
  \centering
  \includegraphics[width=0.88\textwidth]{11_rul_target_shape.png}
  \caption{Raw RUL (dashed) vs capped RUL (solid, cap = 125).
  Early in engine life the signal carries little degradation
  information; the cap replaces uninformative large values
  with a flat healthy label.}
  \label{fig:rulshape}
\end{figure}

\textbf{Why 125?} This is a convention in the research literature,
not a physical constant. Values from 100 to 150 appear in different
papers. The reasoning: sensor signals only start showing degradation
trend roughly 100--150 cycles before failure. Asking a model to
predict RUL = 300 from a healthy-looking sensor trace is both
unreasonable and wastes model capacity.

\textbf{Spreadsheet formula:}
\begin{verbatim}
=MIN(AD2, 125)
\end{verbatim}
(where AD is the RUL\_raw column)

\subsection{cycle\_pct}

Fractional position in the engine's life, normalised to $[0, 1]$:

\begin{equation}
  \texttt{cycle\_pct}(i) = \frac{t_i}{t_{\max}(u_i)}
  \label{eq:cyclepct}
\end{equation}

\textbf{Physical meaning:}
\texttt{cycle\_pct} = 0 at the first cycle of an engine and
\texttt{cycle\_pct} = 1 at failure. Unlike \texttt{time\_cycles},
it is comparable across engines of different total lifespans.
Useful for questions like ``at 50\% of engine life, what do
the sensors look like?''

\textbf{Spreadsheet formula (LO Calc):}
\begin{verbatim}
{=B2/MAX(IF($A:$A=$A2,$B:$B))}
Ctrl+Shift+Enter (array formula)
\end{verbatim}

% -------------------------------------------------------------------
\section{RUL Sheet Columns}
% -------------------------------------------------------------------

The \texttt{RUL} sheet has four columns. The first two
(\texttt{unit\_number}, \texttt{RUL}) come directly from the
NASA \texttt{RUL\_FDxxx.txt} file. The last two are derived.

\subsection{last\_cycle}

The cycle number at which the test trajectory was cut:

\begin{equation}
  \texttt{last\_cycle}(u) = \max\{t_j : u_j = u,\; j \in \text{test}\}
  \label{eq:lastcycle}
\end{equation}

where $u$ is the engine's \texttt{unit\_number}.

\textbf{Where it comes from:} this is the maximum \texttt{time\_cycles}
value for that engine in the \textbf{test} file. It tells you how
far into the engine's life the test trajectory was observed.

\textbf{Not stated in the paper.} This is a logical consequence of
the problem definition: the test file contains cycles 1 through
$t^*$, so $t^* =$ \texttt{last\_cycle}.

\textbf{Spreadsheet formula (LO Calc):}
\begin{verbatim}
{=MAX(IF(test.A:A=A2,test.B:B))}
Ctrl+Shift+Enter (array formula)
\end{verbatim}

\subsection{total\_life\_est}

Estimated total engine life, combining observed and remaining:

\begin{equation}
  \texttt{total\_life\_est}(u) =
  \texttt{last\_cycle}(u) + \texttt{RUL}(u)
  \label{eq:totallife}
\end{equation}

\textbf{Physical meaning:} the engine was observed for
\texttt{last\_cycle} cycles in the test set. It had
\texttt{RUL} cycles remaining. Therefore its estimated
total lifespan is their sum.

\textbf{Not stated in the paper.} Derived for interpretability.
Allows questions like ``what is the distribution of engine lifespans
across the test set?''

\textbf{Spreadsheet formula:}
\begin{verbatim}
=C2+B2
\end{verbatim}
(plain formula, no array entry needed)

% -------------------------------------------------------------------
\section{Equation Reference}
% -------------------------------------------------------------------

\begin{table}[H]
\centering
\caption{All derived column equations with variable sources.}
\label{tab:equations}
\small
\begin{tabular}{@{}lp{5.5cm}p{5.5cm}@{}}
\toprule
Column & Equation & Variables from \\
\midrule
\texttt{dT\_compressor} &
  $T_{30} - T_{24}$ &
  T30 = col 8, T24 = col 7 \\[4pt]

\texttt{dT\_turbine} &
  $T_{50} - T_{30}$ &
  T50 = col 9, T30 = col 8 \\[4pt]

\texttt{condition\_id} &
  $c(i) = \arg\min_k \|x_i - \mu_k\|^2$ &
  $x_i$ = cols 3,4,5; $\mu_k$ from KMeans fit on train \\[4pt]

\texttt{RUL\_raw} &
  $t_{\max}(u_i) - t_i$ &
  $t_i$ = col 2, $t_{\max}$ = max cycle per engine in train \\[4pt]

\texttt{RUL\_capped} &
  $\min(\texttt{RUL\_raw},\; 125)$ &
  RUL\_raw (derived above) \\[4pt]

\texttt{cycle\_pct} &
  $t_i \;/\; t_{\max}(u_i)$ &
  $t_i$ = col 2, $t_{\max}$ = max cycle per engine in train \\[4pt]

\texttt{last\_cycle} &
  $\max\{t_j : u_j = u\}$ over test &
  col 2 of test file, filtered by unit\_number \\[4pt]

\texttt{total\_life\_est} &
  $\texttt{last\_cycle} + \texttt{RUL}$ &
  last\_cycle (derived), RUL = NASA file \\
\bottomrule
\end{tabular}
\end{table}

\begin{keybox}
\textbf{Spreadsheet vs Python boundary.}
Five of the eight derived columns can be computed as spreadsheet
formulas: \texttt{dT\_compressor}, \texttt{dT\_turbine},
\texttt{RUL\_raw}, \texttt{RUL\_capped}, \texttt{cycle\_pct},
\texttt{last\_cycle}, \texttt{total\_life\_est}.
Only \texttt{condition\_id} requires Python, because KMeans is an
iterative optimisation algorithm that cannot be expressed as a
cell formula.
\end{keybox}

% -------------------------------------------------------------------
\end{document}
